\chapter{Perimeter, Area, Volume}\label{ch:g6:measure}

How long is the fence, how big is the field, how much water fits in the
tank? Three different questions, three different quantities ---
perimeter, area, volume --- each with its own units. Confusing them is
the most common mistake in geometry; this chapter sorts them out for
good.

\section{Lengths and perimeter}

\begin{definition}[Perimeter]\label{def:g6:measure:perimeter}
The \emph{perimeter}\index{perimeter} of a figure is the total length
of its border. It is measured in units of length: millimeters (mm),
centimeters (cm), meters (m), kilometers (km), with
\[
1 \text{ km} = 1000 \text{ m}, \qquad
1 \text{ m} = 100 \text{ cm}, \qquad
1 \text{ cm} = 10 \text{ mm}.
\]
\end{definition}

\begin{example}\label{ex:g6:measure:perimeter}
A rectangle of length $L$ and width $w$ has perimeter
\[
P = L + w + L + w = 2 \times (L + w).
\]
For a $7$ cm by $4$ cm rectangle: $P = 2 \times (7 + 4) = 2 \times 11 =
22$ cm. A square of side $c$ has perimeter $4c$.
\end{example}

\begin{proposition}[Circumference of a circle]\label{prop:g6:measure:circle}
The perimeter (or \emph{circumference}) of a circle of radius $r$ is
\[
P = 2 \pi r ,
\]
where $\pi \approx 3.14$ is the same number for every circle.
\end{proposition}
\admitted

\begin{example}
A circular pond has radius $5$ m. Its border measures
$2 \times \pi \times 5 = 10\pi \approx 31.4$ m. Keep the exact value
$10\pi$ as long as possible; round only at the end.
\end{example}

\section{Areas}

\begin{definition}[Area]\label{def:g6:measure:area}
The \emph{area}\index{area} of a figure measures the surface it covers:
how many unit squares fit inside. Units: cm$^2$ (a square of side
$1$ cm), m$^2$, km$^2$, \dots\ Careful:
\[
1 \text{ m}^2 = 100 \times 100 \text{ cm}^2 = 10\,000 \text{ cm}^2
\]
--- one square meter is a $100$ cm by $100$ cm square, so each step of
the units ladder is worth $100$, not $10$.
\end{definition}

\begin{omfigure}
\begin{tikzpicture}[scale=0.68]
  \draw[gray!50, thin] (0,0) grid (7,4);
  \draw[very thick, omDef] (0,0) rectangle (7,4);
  \fill[omThm!25] (0,3) rectangle (1,4);
  \draw[omThm, thick] (0,3) rectangle (1,4);
  \node[font=\small, omThm] at (2.6,3.5) {$1$ unit square};
  \node[font=\small] at (3.5,-0.5) {$7 \times 4 = 28$ unit squares};
\end{tikzpicture}

{\small The area of a rectangle: $7$ columns of $4$ unit squares each,
so $7 \times 4 = 28$ squares in total. This is why area $=$ length
$\times$ width.}
\end{omfigure}

\begin{proposition}[Basic area formulas]\label{prop:g6:measure:formulas}
\begin{center}
\renewcommand{\arraystretch}{1.5}
\begin{tabular}{l|l}
figure & area \\
\hline
rectangle ($L \times w$) & $A = L \times w$ \\
square (side $c$) & $A = c^2$ \\[4pt]
right triangle (legs $a$, $b$) & $A = \dfrac{a \times b}{2}$ \\[4pt]
disk (radius $r$) & $A = \pi r^2$ \\
\end{tabular}
\end{center}
\end{proposition}

\begin{proof}[Proof for the right triangle]
Two copies of a right triangle with legs $a$ and $b$, glued along the
hypotenuse, form an $a \times b$ rectangle. So the triangle's area is
half the rectangle's: $\frac{ab}{2}$. (The rectangle formula is the
unit-square counting above; the disk formula is admitted, see
\cref{ch:g7:areas}.)
\end{proof}

\begin{omfigure}
\begin{tikzpicture}[scale=0.75]
  \draw[thick, fill=omDef!15] (0,0) -- (4,0) -- (0,2.5) -- cycle;
  \draw[thick, omExo, dashed] (4,0) -- (4,2.5) -- (0,2.5);
  \draw[thin] (0.3,0) -- (0.3,0.3) -- (0,0.3);
  \node[below, font=\small] at (2,-0.1) {$a$};
  \node[left, font=\small] at (0,1.25) {$b$};
  \node[font=\small, omExo] at (2.8,1.7) {second copy};
\end{tikzpicture}

{\small A right triangle is half a rectangle: hence the formula
$\frac{a\times b}{2}$.}
\end{omfigure}

\begin{example}[Composite figures]\label{ex:g6:measure:composite}
An L-shaped room is a $6$ m $\times$ $4$ m rectangle with a $2$ m
$\times$ $2$ m square corner removed. Its area, step by step:
\begin{enumerate}
  \item full rectangle: $6 \times 4 = 24$ m$^2$;
  \item removed square: $2 \times 2 = 4$ m$^2$;
  \item remaining area: $24 - 4 = 20$ m$^2$.
\end{enumerate}
Its \emph{perimeter} is not $20$ anything: walking around the L, the
border still measures $6 + 4 + 6 + 4 = 20$ m --- the two cuts of the
corner replace two equal pieces of wall. Same number by coincidence,
but square meters for one, meters for the other!
\end{example}

\begin{remark}[Same perimeter, different areas]
Two figures can have the same perimeter and very different areas: a
$5 \times 5$ square and a $9 \times 1$ rectangle both have perimeter
$20$, but areas $25$ and $9$. Perimeter and area are truly independent
quantities.
\end{remark}

\section{Volumes}

\begin{definition}[Volume]\label{def:g6:measure:volume}
The \emph{volume}\index{volume} of a solid measures the space it fills:
how many unit cubes fit inside. Units: cm$^3$, m$^3$, \dots\ with
$1$ m$^3 = 1\,000\,000$ cm$^3$ (each step of the ladder is worth
$1000$). For liquids one also uses the \emph{liter}:
\[
1 \text{ L} = 1 \text{ dm}^3 = 1000 \text{ cm}^3,
\qquad
1 \text{ m}^3 = 1000 \text{ L}.
\]
\end{definition}

\begin{proposition}[Volume of a box]\label{prop:g6:measure:box}
A rectangular box (a \emph{rectangular prism}) of length $L$, width $w$
and height $h$ has volume
\[
V = L \times w \times h .
\]
\end{proposition}

\begin{proof}
The bottom layer contains $L \times w$ unit cubes
(\cref{def:g6:measure:area} picture, with cubes), and there are $h$
such layers.
\end{proof}

\begin{omfigure}
\begin{tikzpicture}[scale=0.6]
  \foreach \y in {0,1} {
    \foreach \x in {0,...,3} {
      \foreach \z in {0,1} {
        \begin{scope}[shift={({\x+0.45*\z},{\y+0.3*\z})}]
          \draw[fill=omDef!12] (0,0) rectangle (1,1);
          \draw[fill=omDef!20] (0,1) -- (0.45,1.3) -- (1.45,1.3) -- (1,1)
            -- cycle;
          \draw[fill=omDef!28] (1,0) -- (1.45,0.3) -- (1.45,1.3) -- (1,1)
            -- cycle;
        \end{scope}
      }
    }
  }
  \node[below, font=\small] at (2.5,-0.3)
    {$4 \times 2 \times 2 = 16$ unit cubes};
\end{tikzpicture}

{\small A box filled with unit cubes: $4 \times 2$ cubes per layer, two
layers.}
\end{omfigure}

\begin{example}\label{ex:g6:measure:aquarium}
An aquarium measures $60$ cm by $30$ cm by $40$ cm (height). Volume:
\[
V = 60 \times 30 \times 40 = 72\,000 \text{ cm}^3 = 72 \text{ dm}^3
= 72 \text{ L}.
\]
Step by step for the conversion: $1000$ cm$^3$ make one liter, and
$72\,000 \div 1000 = 72$.
\end{example}

\section{Exercises}

\begin{exercise}[$\star$]\label{exo:g6:measure:1}
Convert: $3.5$ m into cm; $420$ mm into cm; $0.8$ km into m;
$25\,000$ m into km.
\end{exercise}

\begin{exercise}[$\star$]\label{exo:g6:measure:2}
Compute the perimeter of: a rectangle $8$ cm $\times$ $3.5$ cm; a
square of side $6.2$ cm; a triangle with sides $5$ cm, $7$ cm and
$9$ cm.
\end{exercise}

\begin{exercise}[$\star$]\label{exo:g6:measure:3}
A circular running track has radius $50$ m. How long is one lap
(exact value with $\pi$, then rounded to the meter)? How many laps make
at least $2$ km?
\end{exercise}

\begin{exercise}[$\star$]\label{exo:g6:measure:4}
Compute the area of: a rectangle $9$ cm $\times$ $4$ cm; a square of
side $7$ m; a right triangle with legs $6$ cm and $10$ cm; a disk of
radius $3$ cm (exact value, then rounded to the cm$^2$).
\end{exercise}

\begin{exercise}[$\star$]\label{exo:g6:measure:5}
Convert: $3$ m$^2$ into cm$^2$; $45\,000$ cm$^2$ into m$^2$; $2.5$ L
into cm$^3$; $4\,500$ L into m$^3$.
\end{exercise}

\begin{exercise}[$\star$]\label{exo:g6:measure:6}
A rectangular field is $120$ m long and $85$ m wide. How many meters of
fence are needed to enclose it? What is its area?
\end{exercise}

\begin{exercise}[$\star$]\label{exo:g6:measure:7}
Compute the volume of a box $5$ cm $\times$ $4$ cm $\times$ $10$ cm,
and of a cube of edge $3$ cm.
\end{exercise}

\begin{exercise}[$\star$]\label{exo:g6:measure:8}
Draw two different rectangles with perimeter $16$ cm, and compute their
areas. Which of your rectangles has the larger area?
\end{exercise}

\begin{exercise}[$\star\star$]\label{exo:g6:measure:9}
A T-shaped figure is made of a $10 \times 2$ horizontal rectangle on
top of a $2 \times 6$ vertical one (measurements in cm). Compute its
area, then its perimeter (walk around the border carefully, adding
every edge).
\end{exercise}

\begin{exercise}[$\star\star$]\label{exo:g6:measure:10}
A swimming pool is a rectangular box $25$ m long, $10$ m wide and
$2$ m deep.
\begin{enumerate}
  \item How many cubic meters of water does it hold when full?
  \item How many liters is that?
  \item The pool is filled at $50\,000$ L per hour. How long does the
        filling take?
\end{enumerate}
\end{exercise}

\begin{exercise}[$\star\star$]\label{exo:g6:measure:11}
A garden is a square of side $20$ m containing a circular pond of
radius $3$ m. What area of grass is there (exact value, then rounded
to the m$^2$)?
\end{exercise}

\begin{exercise}[$\star\star\star$]\label{exo:g6:measure:12}
A chocolate bar measures $15$ cm $\times$ $6$ cm $\times$ $1$ cm.
The maker doubles \emph{all three} dimensions.
\begin{enumerate}
  \item By how much is the volume multiplied? Verify by computing both
        volumes.
  \item The price is multiplied by $4$. Is the big bar a better deal
        than the small one?
\end{enumerate}
\end{exercise}
